\documentclass[screen,oneside]{rdproject}

\begin{filecontents*}[overwrite]{references.bib}
@standard{iso1926,
  title     = {Rigid Cellular Plastics: Determination of Tensile Properties},
  number    = {ISO 1926:2009},
  publisher = {International Organization for Standardization},
  location  = {Geneva, Switzerland},
  date      = {2009}
}
\end{filecontents*}

\addbibresource{references.bib}

\title{Engineering Project Report}
\author{Example Student}
\studentid{202600001}
\addauthor{Second Student}{202600002}

\documentsetup{
  course={Research and Development Project},
  course-code={MECH-501},
  course-short={R\&D Project},
  running-title={Engineering Project Report},
  semester={Autumn 2026},
  supervisor={Example Supervisor},
  submission-date={8 October 2026},
  submission-date-iso={2026-10-08},
  logo={../../figures/AU.pdf}
}

\newglossaryentry{E}{
  name={\ensuremath{E}},
  description={Young's modulus},
  unit={\si{\pascal}},
  sort={A-E}
}
\newglossaryentry{rho}{
  name={\ensuremath{\rho}},
  description={Density},
  unit={\si{\kilogram\per\meter\cubed}},
  sort={B-rho}
}
\newglossaryentry{fem}{
  name={FEM},
  description={Finite element method},
  unit={},
  sort={D-FEM}
}

\begin{document}

\frontmatter
\maketitle

\begin{abstract}
This example demonstrates the public structure and formatting provided by the
\texttt{rdproject} class.
\end{abstract}

\begin{preface}
The report class provides project metadata, front matter, nomenclature,
IEEE-style references, appendices, and project-specific running headers.
\end{preface}

\tableofcontents
\glsaddall
\printglossary[style=NoahNomenclature,title={Nomenclature}]
\printlistoffigures
\printlistoftables

\mainmatter

\chapter{Introduction}
The project follows the tensile-testing terminology of \cite{iso1926}.
The symbol \gls{E} denotes Young's modulus and \gls{rho} denotes density.

\section{Example equation}
A simple constitutive relation is
\begin{equation}
  \sigma = E\varepsilon.
  \label{eq:hooke}
\end{equation}
As shown in \Mref{eq:hooke}, stress is proportional to strain in this idealised
linear-elastic example.

\section{Example figure and table}
\begin{figure}[htbp]
  \centering
  \fbox{\rule{0pt}{35mm}\rule{0.65\linewidth}{0pt}}
  \caption{Placeholder project figure.}
  \label{fig:placeholder}
\end{figure}

\begin{table}[htbp]
  \centering
  \caption{Example material properties.}
  \label{tab:properties}
  \begin{tabular}{lcc}
    \toprule
    Material & \(E\) & \(\rho\) \\
    \midrule
    Example A & \qty{2.0}{\giga\pascal} & \qty{100}{\kilogram\per\meter\cubed} \\
    Example B & \qty{3.0}{\giga\pascal} & \qty{120}{\kilogram\per\meter\cubed} \\
    \bottomrule
  \end{tabular}
\end{table}

\appendix
\chapter{Additional Derivation}\label{app:derivation}
This appendix demonstrates appendix-aware cross-references.

\backmatter
\printreferences

\end{document}
